On se place dans le plan muni d'un repère orthonormé.

On considère les points $A(-1\ ;\ 5)$, $B(3\ ;\ 5)$ et $C(4\ ;\ 0)$.

\begin{enumerate}
\item \begin{enumerate}
\item Calculer les coordonnées des vecteurs $\overrightarrow{AB}$ et $\overrightarrow{AC}$.
\item En déduire la valeur du produit scalaire $\overrightarrow{AB}\cdot\overrightarrow{AC}$.
\end{enumerate}
\item \begin{enumerate}
\item Montrer que $AC=5\sqrt{2}$.

On admet que $AB=4$.
\item Écrire l'expression permettant de calculer le produit scalaire $\overrightarrow{AB}\cdot\overrightarrow{AC}$ en fonction de l'angle $\widehat{BAC}$.
\item En déduire une mesure, en radian, de l'angle $\widehat{BAC}$.
\end{enumerate}
\end{enumerate}

\vfill
\begin{center}
\begin{tblr}{colspec={Q[l]},width=0.55\linewidth}
\textbf{Aide aux calculs :} \\
$\dfrac{1}{\sqrt{2}}=\dfrac{\sqrt{2}}{2}$ \\
$\dfrac{5}{\sqrt{50}}=\dfrac{\sqrt{2}}{2}$ \\
\end{tblr}
\end{center}

\clearpage
